\clearpage
\section*{Ответы для самопроверки}
Сначала решите задачи самостоятельно. Номера соответствуют каждому практикуму.
\subsection*{Практикум 1}
\begin{multicols}{3}
\begin{enumerate}
\item $\dfrac25$.
\item $\dfrac14$.
\item $6$.
\item $\dfrac34$.
\item $-\dfrac13$.
\item $-\dfrac43$.
\item $-\dfrac12$.
\item $\text{не существует}$.
\item $\dfrac13$.
\item $2$.
\item $-\dfrac13$.
\item $\dfrac1{48}$.
\item $\dfrac12$.
\item $4$.
\item $-\dfrac5{12}$.
\item $0$.
\item $-\sqrt5$.
\item $\dfrac12$.
\item $\dfrac23$.
\item $0$.
\item $+\infty$.
\item $\dfrac14$.
\end{enumerate}
\end{multicols}
В задаче 8 предел слева равен $+\infty$, справа --- $-\infty$. В задаче 16 подход к нулю возможен только справа. В задаче 21 кубический корень понимается как действительный.
\subsection*{Практикум 2}
\begin{multicols}{3}
\begin{enumerate}
\item $\dfrac23$.
\item $5$.
\item $\dfrac32$.
\item $\dfrac1{20}$.
\item $6$.
\item $-3$.
\item $\dfrac{25}{4}$.
\item $8$.
\item $5$.
\item $\dfrac12$.
\item $2$.
\item $4$.
\item $2$.
\item $1$.
\item $\dfrac12$.
\item $\dfrac13$.
\item $1$.
\item $0$.
\item $-2$.
\item $\dfrac9{98}$.
\item $1$.
\end{enumerate}
\end{multicols}
\subsection*{Практикум 3}
\begin{multicols}{3}
\begin{enumerate}
\item $+\infty$.
\item $\dfrac35$.
\item $-2$.
\item $3$.
\item $1$.
\item $e$.
\item $e^{-7}$.
\item $e$.
\end{enumerate}
\end{multicols}
